\chapter{Trabalhos Relacionados}
\label{cap:trabalhos-relacionados}
'
A aplicação de técnicas de aprendizado profundo à classificação automatizada do câncer de próstata tem sido amplamente investigada na literatura, com avanços expressivos particularmente no uso de imagens de lâminas histopatológicas completas (Whole-Slide Images – WSIs), estratégias baseadas na extração de patches e mecanismos de agregação em múltiplas escalas. Esta seção apresenta uma revisão dos principais trabalhos relacionados, organizados de acordo com a metodologia predominante adotada.

\section{Aprendizagem Fracamente Supervisionada}
\cite{Xiang2023Prostate} propuseram o framework GCN-MIL, que utiliza uma Rede Neural de Grafos para agregar características espaciais de patches em lâminas de WSI. O modelo incorpora a estratégia de \textit{Robust Training}, que filtra iterativamente amostras de alta incerteza para mitigar o impacto de rótulos ruidosos e inconsistentes, permitindo alcançar um Kappa quadrático de 0,931 no conjunto interno do desafio PANDA. Entretanto, a abordagem apresenta lacunas fundamentais ao utilizar a função de perda \textit{Cross-Entropy} tradicional, que trata as graduações ISUP como classes nominais independentes e ignora a hierarquia biológica da progressão do câncer de próstata. Além disso, a dependência de uma arquitetura ResNet50 para a extração de características limita o potencial de representação morfológica e a eficiência computacional em relação a arquiteturas mais modernas, como a EfficientNet.

\cite{jung2022artificial} realizaram uma validação externa independente do sistema DeepDx Prostate utilizando 593 imagens de lâmina inteira do Seoul National University Hospital , alcançando um coeficiente kappa ponderado quadrático de 0,922 na graduação de Gleason. O sistema demonstrou desempenho superior aos laudos originais na detecção de padrões de Gleason 4 e 5, elevando a concordância diagnóstica de um patologista geral de 0,876 para 0,925. Apesar dos bons resultados, a generalização é limitada pela homogeneidade da amostra, restrita a uma única instituição e a um único modelo de scanner , além de uma tendência observada de superestimação de casos do Grupo de Grau 1 (GG1), em que a alta sensibilidade do algoritmo a pequenas frações dos padrões 4 ou 5 resultou na reclassificação inadequada de biópsias de baixo grau para categorias de maior gravidade.

\cite{liu2022deep} investigaram o uso de aprendizado profundo para identificar pacientes com risco de câncer de próstata a partir de lâminas histopatológicas classificadas como benignas. Os autores empregaram um pipeline baseado em WSIs, com segmentação de tecido, extração massiva de patches e um ensemble de redes convolucionais treinadas sob supervisão fraca, no qual as predições em nível de \textit{patch} foram agregadas hierarquicamente para inferência em nível de lâmina e de paciente. Avaliado em um conjunto independente, o método alcançou AUC de 0,739 na predição de risco em nível de paciente, superando abordagens clínicas baseadas apenas na idade e no PSA. Análises interpretáveis com UMAP e Grad-CAM indicaram que o modelo capturou padrões morfológicos sutis em tecido benigno associados à presença de câncer não amostrado, especialmente características glandulares e nucleares atípicas. Apesar do potencial clínico demonstrado para redução de falsos negativos e apoio à decisão de rebiopsia, o estudo apresenta limitações relacionadas à dependência de um único domínio institucional e à necessidade de validação adicional para garantir robustez e generalização interlaboratorial.

\cite{pyramidalcnn} propuseram um sistema automatizado para graduação do câncer de próstata em imagens histopatológicas digitalizadas de biópsias utilizando uma arquitetura Pyramidal CNN composta por três redes convolucionais rasas operando em diferentes escalas de patch (100×100, 150×150 e 200×200 pixels), permitindo capturar características morfológicas em múltiplos níveis de contexto espacial. O pipeline de processamento incluiu equalização de histograma para normalização de intensidade, extração de patches a partir das WSIs e seleção com base na proporção mínima de tecido presente, seguida de classificação patch-wise dos padrões de Gleason (GP1–GP5) e agregação das probabilidades para inferência do Gleason Score e dos Grade Groups. O conjunto experimental foi composto por 896 lâminas inteiras de biópsias prostáticas digitalizadas (750 para treinamento, 96 para validação e 50 para teste), fornecidas pelo Radboud University Medical Center e analisadas no University of Louisville Hospital, com rotulagem semi-automática realizada por patologistas especialistas. Os experimentos geraram milhões de \textit{patches} multiescala para treinamento das redes, e os resultados indicaram acurácia de aproximadamente 0,77 na classificação dos padrões, além de desempenho na identificação dos grupos prognósticos com recall entre 50\% e 75\% e precisão entre 50\% e 69\%, comparável ao desempenho humano em cenários similares. Contudo, o método apresenta limitações relacionadas ao uso de rotulação semiautomática, à ausência de validação externa independente, ao desempenho moderado em classes de maior complexidade e à dependência da estratégia de seleção de \textit{patches}, fatores que restringem a avaliação de generalização clínica e robustez em ambientes multi-institucionais.

No trabalho de \cite{kosoko2024xai_prostate}, os autores avaliaram diferentes arquiteturas de redes neurais convolucionais, combinadas com técnicas de \textit{Explainable Artificial Intelligence} (XAI), em especial \textit{Grad-CAM}, para a interpretação das decisões do modelo. O conjunto de dados utilizado foi o SICAPv2, composto por patches extraídos de WSIs coradas em H\&E com anotação em nível de pixel, incluindo 4417 amostras não cancerosas, 1635 do grupo GG3, 3622 do GG4 e 665 do GG5. Entre as arquiteturas avaliadas, o modelo VGG19 apresentou o melhor desempenho, alcançando AUC-ROC de 0{,}85 e valores médios próximos de 0{,}75 para acurácia, precisão e \textit{recall}, enquanto uma CNN simples apresentou desempenho significativamente inferior (acurácia $\approx$ 0{,}41 e F1 $\approx$ 0{,}24), evidenciando a vantagem de arquiteturas profundas pré-treinadas na captura de padrões morfológicos complexos associados aos graus de Gleason. Apesar dos resultados promissores, o estudo apresenta limitações relacionadas à ausência de validação clínica externa, à dependência de um único conjunto de dados e à exploração restrita de técnicas de XAI, o que limita a avaliação de generalização em cenários interinstitucionais.

No estudo de \cite{afifi2024prostate}, a detecção de câncer de próstata em imagens histopatológicas digitalizadas foi formulada como um problema de classificação binária, visando distinguir tecidos benignos de malignos. Para essa finalidade, foram avaliadas arquiteturas profundas pré-treinadas, InceptionV3, ResNet50 e InceptionResNetV2, posteriormente ajustadas por meio de \textit{fine-tuning} para a extração de características morfológicas relevantes. O conjunto de dados utilizado é composto por imagens de lâminas histopatológicas provenientes da Radboud University Medical Center e da Karolinska Institute, organizadas em formatos WSI e convertidas em patches para treinamento supervisionado. As imagens foram redimensionadas, normalizadas e submetidas a técnicas de aumento de dados, e posteriormente divididas em subconjuntos de treinamento, validação e teste. O treinamento utilizou a função de perda \textit{binary cross-entropy}, evidenciando a modelagem do problema como detecção de câncer e não como classificação multiclasse de padrões de Gleason. Os resultados experimentais indicaram desempenho superior da arquitetura InceptionResNetV2, que atingiu acurácia de 91,8\%, demonstrando elevada capacidade de generalização na distinção entre tecidos saudáveis e cancerosos.

\section{Atenção e Transformers}

\cite{alici2024efficient} propuseram um modelo de aprendizado profundo para diagnóstico e graduação do câncer de próstata utilizando imagens histopatológicas do dataset público DiagSet. O processamento foi realizado em patches extraídos de lâminas digitalizadas coradas em H\&E. A abordagem baseou-se na arquitetura EfficientNet-B4 e em uma variante denominada Eff4-Attn, que incorpora um mecanismo de atenção de canais para aprimorar a extração de características morfológicas. O modelo foi avaliado em diferentes níveis de granularidade: uma classificação em quatro classes (tecido normal, estroma, alterações benignas e câncer) e outra em oito classes (tecido normal, estroma, artefatos, inflamação, hiperplasia benigna, Gleason 3, Gleason 4 e Gleason 5). Os resultados alcançaram acurácia de 96,18\% na tarefa binária, 94,86\% na classificação em quatro classes e 93,32\% na classificação em oito classes, demonstrando desempenho competitivo aliado à eficiência computacional. Contudo, o estudo apresenta limitações relacionadas ao uso exclusivo de patches em vez de inferência em nível de lâmina completa, à dependência de um único conjunto de dados público e ao desbalanceamento entre as classes, fatores que restringem a avaliação da generalização clínica e da robustez do método em cenários multi-institucionais.

No estudo de \cite{ikromjanov2021vit_prostate}, a análise de imagens histopatológicas de próstata foi conduzida por meio de um modelo baseado em \textit{Vision Transformers} (ViT), com o objetivo de classificar padrões teciduais associados ao sistema de graduação de Gleason a partir de imagens WSI. O conjunto de dados utilizado deriva do desafio público PANDA, sendo empregado um subconjunto com mais de 5.000 imagens na experimentação. Após a detecção das regiões de interesse, as lâminas foram particionadas em aproximadamente 304 mil \textit{patches} de dimensões 256×256 pixels, rotuladas em cinco classes (estroma, benigno, Gleason 3, Gleason 4 e Gleason 5), com cerca de 240 mil amostras destinadas ao treinamento, 60 mil à validação e 4 mil ao teste. O modelo ViT foi utilizado para classificação em nível de \textit{patch} e para inferência em nível de lâmina, com predições agregadas. Os resultados experimentais demonstraram desempenho competitivo na identificação de padrões histológicos, alcançando acurácia superior a 85\% na classificação dos \textit{patches}, evidenciando a capacidade do mecanismo de atenção de capturar dependências contextuais globais nas imagens. 

\section{Limitações da Literatura e Lacunas de Pesquisa}

Apesar dos avanços na aplicação de aprendizado profundo à análise de imagens histopatológicas de próstata, a literatura ainda apresenta limitações metodológicas relevantes. Diversos trabalhos formulam o problema como uma classificação nominal e utilizam funções de perda tradicionais, como \textit{Cross-Entropy} e \textit{Binary Cross-Entropy}, desconsiderando a natureza ordinal dos graus ISUP. Como consequência, erros de predição são penalizados de forma uniforme, independentemente da magnitude da divergência entre as classes, o que não reflete adequadamente a relevância clínica das discrepâncias na graduação tumoral. Além disso, muitas abordagens dependem exclusivamente de arquiteturas convolucionais com extração local de características ou de estratégias de agregação baseadas em \textit{patches}, sem incorporar mecanismos robustos de filtragem de ruído ou de modelagem de incerteza. Tal limitação é particularmente crítica diante da presença conhecida de artefatos de coloração, inconsistências de anotação e variabilidade interinstitucional em conjuntos de dados públicos amplamente utilizados, como o PANDA, o que pode comprometer a capacidade de generalização dos modelos.

Outra limitação recorrente refere-se à dependência de classificadores individuais, com pouca exploração sistemática de ensembles heterogêneos capazes de aumentar a robustez e a estabilidade preditiva. Da mesma forma, abordagens que utilizam classificação binária ou multiclasses plana frequentemente deixam de capturar relações hierárquicas entre graus tumorais, reduzindo o alinhamento clínico dos modelos com o processo diagnóstico real.

Diante dessas lacunas, o presente trabalho propõe um pipeline estruturado que integra a filtragem baseada em entropia para redução de ruído, a modelagem ordinal com Focal Loss para a captura da progressão biológica da doença e estratégias de ensemble entre variantes do EfficientNet, buscando melhorar simultaneamente a robustez, a estabilidade e a coerência clínica das predições na classificação ISUP.
